\documentclass[10pt,letterpaper]{article}
\usepackage{cloudsuite}
\renewcommand{\DocShort}{CH. 03}
\begin{document}
\coverpage{Chapter 3: Understanding Cloud Computing}{Cornell study notes and applied review | FOUNDATIONS}{Concepts \textbullet{} decisions \textbullet{} verification artifacts}
\sourceStatement
\tableofcontents
\newpage
\section{Learning contract}
\begin{takeaway}
Reason about business value and exposure before selecting an architecture.
\end{takeaway}
\textbf{Chapter scope:} Origins, terminology, benefits and adoption risks.\\
\sourceRef{3}{23}
\subsection{Mapped source sections}
\begin{itemize}
\item 3.1 Origins and Influences
\item 3.2 Basic Concepts and Terminology
\item 3.3 Goals and Benefits
\item 3.4 Risks and Challenges
\end{itemize}
\section{Cornell cue and notes matrix}
\small\setlength{\LTleft}{0pt}\setlength{\LTright}{0pt}
\begin{longtable}{>{\RaggedRight\arraybackslash}p{0.29\textwidth}>{\RaggedRight\arraybackslash}p{0.63\textwidth}}
\toprule\textbf{Cue / essential question}&\textbf{Explanatory notes and reasoning}\\\midrule\endfirsthead\toprule\textbf{Cue}&\textbf{Notes}\\\midrule\endhead
\textbf{How did modern cloud practices emerge?} & Clustering, grid-oriented resource sharing, virtualization, containers and serverless environments represent different contributors to elastic, remotely provisioned computing. Do not treat them as synonymous.\\[3pt]\textcolor{petrol}{\emph{Practice:} Explain what is abstracted at each layer.}\\[8pt]\midrule
\textbf{What is the provider-consumer boundary?} & A cloud provider operates the environment that offers IT resources; cloud consumers obtain and use services under agreed conditions. Service owners and administrators can be different actors.\\[3pt]\textcolor{petrol}{\emph{Practice:} Mark ownership and access on a simple system diagram.}\\[8pt]\midrule
\textbf{How do horizontal and vertical scaling differ?} & Horizontal scaling changes instance count; vertical scaling changes capacity within an instance. Both depend on bottlenecks, state and resource ceilings.\\[3pt]\textcolor{petrol}{\emph{Practice:} Predict which scaling direction fits a stateful database and a stateless API.}\\[8pt]\midrule
\textbf{Why is elastic consumption economically attractive?} & Provisioning to demand can reduce unused reserved capacity, but usage charging shifts control toward instrumentation and budgets.\\[3pt]\textcolor{petrol}{\emph{Practice:} Estimate a burst-month cost with an explicit unit rate.}\\[8pt]\midrule
\textbf{What are the recurring adoption risks?} & Shared responsibility, overlapping trust boundaries, portability limits, operational control gaps, legal exposure and cost overrun can offset anticipated benefits.\\[3pt]\textcolor{petrol}{\emph{Practice:} Create a risk register with accountable control owners.}\\[8pt]\midrule
\textbf{How can availability claims be misunderstood?} & More resources do not automatically eliminate common-mode failures. Failure domains and continuity tests must accompany a promise of reliability.\\[3pt]\textcolor{petrol}{\emph{Practice:} Draw two allegedly redundant systems that actually share a failure point.}\\[8pt]\midrule
\end{longtable}\normalsize
\section{Synthesis and architecture connections}
Translate the chapter vocabulary into one explicit service requirement, one boundary, one owner and one measurable test. Later mechanism and architecture chapters provide possible realizations; the requirements remain the acceptance baseline.
\begin{workbox}[title={Applied deliverable}]
Reason about business value and exposure before selecting an architecture.\\[4pt]Create a one-page architecture or operations artifact with: (1) assumptions, (2) a decision and alternatives, (3) control ownership, (4) observable acceptance criteria, and (5) a failed-case test.
\end{workbox}
\section{Active recall}
\begin{enumerate}
\item How did modern cloud practices emerge?
\item What is the provider-consumer boundary?
\item How do horizontal and vertical scaling differ?
\item Why is elastic consumption economically attractive?
\end{enumerate}
\subsection{Answer-check rubric}
\begin{itemize}
\item A correct response distinguishes the named concepts rather than relying on vendor names.
\item The explanation identifies an operational consequence, boundary or tradeoff.
\item At least one measurement, test or documentary artifact could demonstrate the claim.
\end{itemize}
\section{Blank Cornell capture sheet}
\textbf{Topic / source subsection:}\enspace\rule{0.65\linewidth}{0.4pt}\\[6pt]
\begin{longtable}{p{0.28\textwidth}p{0.64\textwidth}}
\toprule Cues / questions & Notes / evidence \\ \midrule
\rule{0pt}{42pt} & \\ \midrule
\rule{0pt}{42pt} & \\ \midrule
\rule{0pt}{42pt} & \\ \midrule
\rule{0pt}{42pt} & \\ \midrule
\rule{0pt}{42pt} & \\ \midrule
\end{longtable}
\textbf{Summary in three sentences:}\\[16pt]\hrulefill\\[16pt]\hrulefill\\[16pt]\hrulefill

\refs
\end{document}
